\documentclass[11pt]{article}
\usepackage[margin=0.9in]{geometry}
\usepackage{amsmath,amssymb,amsthm,booktabs,mathrsfs}
\usepackage[T1]{fontenc}
\usepackage{lmodern,microtype}
\usepackage[colorlinks=true,linkcolor=blue,urlcolor=blue]{hyperref}
\newtheorem{theorem}{Theorem}[section]
\newtheorem{lemma}[theorem]{Lemma}
\newtheorem{proposition}[theorem]{Proposition}
\newtheorem{corollary}[theorem]{Corollary}
\theoremstyle{remark}
\newtheorem{remark}[theorem]{Remark}
\newcommand{\norm}[1]{\left\|#1\right\|}
\title{Toward a four-dimensional Yang--Mills construction\\
\large Notebook 30: quantum descent and renormalized block stability}
\author{Research working document}
\date{17 October 2026}
\begin{document}
\maketitle
\begin{abstract}
We continue the existence argument by constructing the local quantum
update on the complete interaction quotient of Notebook 29.
Actual finite minimum comparisons supply connecting gauges that,
in the infinite local limit, approach one constant exponentially.
Compact gauge approximation then proves that classical substitution
preserves local gauge invariance for every invariant input class.
A norm-preserving periodic construction identifies finite normalized
interactions independently of their carrier representation, including
wrapping supports and exact finite Ward identities.

The actual strict Gaussian expansion is continuous in its local
analytic data. Its proof uses support incidence, strict marked-order
tails, an explicit positive inherited-cardinality margin and finite
dimensional dominated convergence. These estimates define the
homogeneous local quantum map on the invariant quotient and bound
the input dependence of its quantum remainder by
\(C(gT/b)\). Finite translation averaging identifies limits of varying
finite inputs with the same map, giving compatibility for every fixed
finite history inside the common budgets.

We then examine the renormalized recurrence using the actual classical
reference actions and the finite endpoint ratio \(R^q\).
Exact reference transport cancels intermediate classical projections.
A fixed outer field domain and vanishing total linear jet give
an \(O(b^2)\) zero-input quantum force on the inner radius \(b\).
The extracted activities therefore have an \(O(b^2)\) invariant
envelope under a finite block of sufficiently weak steps, with
all intermediate inputs in the common domain. Every quantum
correction and vacuum factor remains present. These estimates concern
the homogeneous local sector relative to one fixed joint classical
system. The running
coupling must obey its actual extraction equation; a prescribed
sequence of small parameters does not by itself furnish that
trajectory. Mixed and large-field histories, sourced bounds and
the continuum observable theory remain necessary for existence.
\end{abstract}

The objective is a nontrivial four-dimensional Yang--Mills quantum
theory. All results here use the finite constructions of Notebooks
15--29 with their stated common domains and actual retained-link
windows. The proofs below add explicit bridges between finite
normalized integrals, infinite analytic interaction classes and
renormalized activity estimates. They do not assume that equality
of infinite compact-field totals automatically identifies a finite
integral, nor that a final norm bound alone proves continuity.

